%%%%%%%%%%%%%%%%%%%%%%%%%%%%%%%%%%%%%%%%%%%%%%%%%%%%%%%%%%%%%%%%%%%%%%%%%%%%%%%%%%
\begin{frame}[fragile]\frametitle{}
\begin{center}
{\Large Case Study: Grouping Cars}
\end{center}
\end{frame}

%%%%%%%%%%%%%%%%%%%%%%%%%%%%%%%%%%%%%%%%%%%%%%%%%%%%%%%%%%
\begin{frame}[fragile]\frametitle{Cars: The Data}
\begin{itemize}
\item 428 cars: engine size, cylinders, horsepower, city MPG, weight, and more
\item Data: the SAS sample data set \texttt{SASHELP.CARS}
\item Goal: find natural vehicle groups; K-Means is given no labels
\item Clean first: the Mazda RX-8 (rotary engine) has \texttt{'.'} for cylinders, so 2 cars are dropped
\end{itemize}
\begin{lstlisting}
import pandas as pd

cars = pd.read_csv('data/cars.csv', sep=';')
cars['Cylinders'] = pd.to_numeric(cars['Cylinders'],
                                  errors='coerce')  # '.' becomes NaN

features = ['EngineSize', 'Cylinders', 'Horsepower',
            'MPG_City', 'Weight']
cars = cars.dropna(subset=features)     # 428 -> 426 cars
X = cars[features]
\end{lstlisting}
\end{frame}

%%%%%%%%%%%%%%%%%%%%%%%%%%%%%%%%%%%%%%%%%%%%%%%%%%%%%%%%%%
\begin{frame}[fragile]\frametitle{Cars: Cluster the Raw Features}
\begin{itemize}
\item Run K-Means with $k=3$ on the raw columns
\item Weight is in the thousands; cylinders are single digits
\item Weight's variance is about 577,000, against at most 5,200 for any other feature
\item Result: the three groups are exactly the groups Weight alone gives
\end{itemize}
\begin{lstlisting}
from sklearn.cluster import KMeans

km = KMeans(n_clusters=3, n_init=10, random_state=0)
km.fit(X)
cars['cluster'] = km.labels_

# In the Python shell (cluster averages):
>>> cars.groupby('cluster')[features].mean().round(1).T
cluster          0       1       2
EngineSize     3.4     2.2     4.6
Cylinders      6.1     4.3     7.5
Horsepower   231.7   154.4   275.8
MPG_City      18.7    25.3    15.0
Weight      3635.4  2781.1  4874.4
\end{lstlisting}
\end{frame}

%%%%%%%%%%%%%%%%%%%%%%%%%%%%%%%%%%%%%%%%%%%%%%%%%%%%%%%%%%
\begin{frame}[fragile]\frametitle{Cars: Scale First}
\begin{itemize}
\item Standardize each feature (mean 0, spread 1): all count equally
\item 107 of 426 cars (25\%) change group versus the raw features
\item Now all five features shape the groups: small economy cars, mid-size 6-cylinder cars, large heavy 8-cylinder cars
\item Cluster numbers are arbitrary: compare the averages, not the numbers (cluster 0 here is not cluster 0 on the previous slide)
\end{itemize}
\begin{lstlisting}
from sklearn.preprocessing import StandardScaler

Xs = StandardScaler().fit_transform(X)
km = KMeans(n_clusters=3, n_init=10, random_state=0)
km.fit(Xs)
cars['cluster'] = km.labels_

# In the Python shell (cluster averages):
>>> cars.groupby('cluster')[features].mean().round(1).T
cluster          0       1       2
EngineSize     2.0     3.2     4.8
Cylinders      4.0     5.9     8.1
Horsepower   144.8   216.7   307.2
MPG_City      25.7    18.7    15.7
Weight      2828.1  3638.1  4445.5
\end{lstlisting}
\end{frame}

%%%%%%%%%%%%%%%%%%%%%%%%%%%%%%%%%%%%%%%%%%%%%%%%%%%%%%%%%%
\begin{frame}[fragile]\frametitle{Cars: Is k=3 Right?}
\begin{itemize}
\item Try $k=2$ to $6$ on the scaled data: inertia and silhouette score
\item Inertia drops sharply up to $k=3$, then only gently: an elbow at $k=3$
\item The silhouette score is highest at $k=3$ (0.457)
\item Both are guides: an engineer may still choose $k=4$
\end{itemize}
\begin{lstlisting}
from sklearn.metrics import silhouette_score

for k in range(2, 7):
    km = KMeans(n_clusters=k, n_init=10, random_state=0)
    km.fit(Xs)
    sil = silhouette_score(Xs, km.labels_)
    print(k, round(km.inertia_, 1), round(sil, 3))

# printed: k, inertia, silhouette
# 2 1079.5 0.409
# 3 653.6 0.457
# 4 548.0 0.429
# 5 468.3 0.412
# 6 402.6 0.384
\end{lstlisting}
\end{frame}

%%%%%%%%%%%%%%%%%%%%%%%%%%%%%%%%%%%%%%%%%%%%%%%%%%%%%%%%%%
\begin{frame}[fragile]\frametitle{Quick Check: Cars and Scaling}
\begin{block}{Think About It}
In the cars data, Weight is recorded in pounds (1,850 to 7,190). A colleague converts it to tonnes (dividing by about 2,200) and re-runs K-Means with $k=3$ on the raw, unscaled columns. What is most likely to happen?
\begin{enumerate}
  \item[A.] Nothing changes: K-Means does not depend on units
  \item[B.] The groups change, because Weight now counts far less in the distance and the other features count more
  \item[C.] Only the Weight column of the cluster centers changes; the groups stay the same
  \item[D.] K-Means fails, because tonnes are not whole numbers
\end{enumerate}
\end{block}
\end{frame}

%%%%%%%%%%%%%%%%%%%%%%%%%%%%%%%%%%%%%%%%%%%%%%%%%%%%%%%%%%%%%%%%%%%%%%%%%%%%%%%%%%
\begin{frame}[fragile]\frametitle{Quick Check: Cars and Scaling}
\begin{block}{Answer}
\textbf{Correct: B.} Distances are computed in raw units, so a unit change silently changes the clustering; on the cars data the groups change markedly. Standardizing first removes the problem: after scaling, pounds and tonnes give identical groups. A and C overlook that the distances themselves change; D is not a real constraint.
\end{block}
\end{frame}
